\frfilename{mt131.tex}
\versiondate{18.3.05}
\copyrightdate{1995}

\def\chaptername{Pelengkap}
\def\sectionname{Subruang terukur}

\newsection{131}

Sangat sering kita ingin mengintegralkan suatu fungsi pada sebuah
subhimpunan dari suatu ruang ukur; misalnya, membentuk integral
$\int_a^bf(x)dx$, dengan $a<b$ dalam $\Bbb R$. Sering kali kita ingin
melakukan ini ketika fungsi tersebut tidak didefinisikan pada sebagian
atau seluruh bagian di luar subhimpunan itu, seperti dalam ungkapan
$\int_0^1\ln x\,dx$. Kerangka alami untuk melakukan operasi semacam itu
adalah `ukuran subruang'. Jika $(X,\Sigma,\mu)$ suatu ruang ukur dan
$H\in\Sigma$, terdapat ukuran subruang alami $\mu_H$ pada $H$, yang
saya uraikan dalam bagian ini.
Saya mulai dengan definisi ukuran subruang ini (131A-131C), disertai
uraian tentang integrasi terhadapnya (131E-131H); ini memberikan
landasan yang kukuh bagi konsep `integrasi pada suatu subhimpunan
(terukur)' (131D).

\leader{131A}{Proposisi} Misalkan $(X,\Sigma,\mu)$ suatu ruang ukur, dan
$H\in\Sigma$. Tetapkan $\Sigma_H=\{E:E\in\Sigma,\,E\subseteq H\}$ dan
misalkan $\mu_H$ merupakan pembatasan $\mu$ pada $\Sigma_H$. Maka
$(H,\Sigma_H,\mu_H)$ merupakan suatu ruang ukur.

\proof{ Tentu saja $\Sigma_H$ tidak lain adalah
$\{E\cap H:E\in\Sigma\}$, dan telah saya catat (dalam 121A) bahwa ini
merupakan suatu aljabar-$\sigma$ atas subhimpunan-subhimpunan $H$.
Sekarang jelas bahwa $\mu_H$ memenuhi (iii) dari 112A, sehingga
$(H,\Sigma_H,\mu_H)$ merupakan suatu ruang ukur.
}%end of proof of 131A

\leader{131B}{Definisi} Jika $(X,\Sigma,\mu)$ sembarang ruang ukur dan
$H\in\Sigma$, maka $\mu_H$, sebagaimana didefinisikan dalam 131A,
adalah {\bf ukuran subruang} pada $H$.

Jika $X=\BbbR^r$, dengan $r\ge 1$, dan $\mu$ adalah ukuran Lebesgue
pada $\BbbR^r$, saya akan menyebut suatu ukuran subruang $\mu_H$
sebagai {\bf ukuran Lebesgue pada} $H$.

\cmmnt{Fakta-fakta elementer berikut patut dicatat.}

\leader{131C}{Lema} Misalkan $(X,\Sigma,\mu)$ suatu ruang ukur,
$H\in\Sigma$, dan $\mu_H$ ukuran subruang pada $H$, dengan domain
$\Sigma_H$. Maka

(a) untuk setiap $A\subseteq H$, $A$ bersifat $\mu_H$-terabaikan jika
dan hanya jika bersifat $\mu$-terabaikan;

(b) jika $G\in\Sigma_H$, maka $(\mu_H)_G$, yaitu ukuran subruang pada
$G$ ketika $G$ dipandang sebagai subhimpunan terukur dari $H$, identik
dengan $\mu_G$, yaitu ukuran subruang pada $G$ ketika $G$ dipandang
sebagai subhimpunan terukur dari $X$.

\leader{131D}{Integrasi pada subhimpunan: Definisi} Misalkan
$(X,\Sigma,\mu)$ suatu ruang ukur, $H\in\Sigma$, dan $f$ suatu fungsi
bernilai real yang didefinisikan pada suatu subhimpunan dari $X$. Dengan
$\int_H f$ (atau $\int_H f(x)\mu(dx)$, dan sebagainya) saya akan
memaksudkan $\int(f\restr H)d\mu_H$, jika integral ini ada\cmmnt{,
sesuai dengan definisi-definisi 131A-131B dan 122M, serta dengan
menetapkan $\dom(f\restr H)=H\cap\dom f$,
$(f\restr H)(x)=f(x)$ untuk $x\in H\cap\dom f$}.


\leader{131E}{Proposisi} Misalkan $(X,\Sigma,\mu)$ suatu ruang ukur,
$H\in\Sigma$, dan $f$ suatu fungsi bernilai real yang didefinisikan
pada suatu subhimpunan $\dom f$ dari $H$. Tetapkan $\tilde f(x)=f(x)$
jika $x\in\dom f$, dan $0$ jika $x\in X\setminus H$. Maka, jika salah
satu ruas dalam $\int fd\mu_H=\int \tilde fd\mu$ terdefinisi dalam
$\Bbb R$, ruas lainnya juga demikian, dan persamaan tersebut berlaku.

\proof{{\bf (a)} Jika $f$ merupakan fungsi $\mu_H$-sederhana, fungsi
itu dapat dinyatakan sebagai $\sum_{i=0}^na_i\chi E_i$, dengan
$E_0,\ldots,E_n\in\Sigma_H$, $a_0,\ldots,a_n\in\Bbb R$, dan
$\mu_H E_i<\infty$ untuk setiap $i$. Sekarang $\tilde f$ juga sama
dengan $\sum_{i=0}^na_i\chi E_i$ jika ungkapan ini ditafsirkan sebagai
fungsi dari $X$ ke $\Bbb R$. Jadi

\Centerline{$\sum_{i=0}^na_i\mu_HE_i
=\sum_{i=0}^na_i\mu E_i=\int\tilde fd\mu$.}

\medskip

{\bf (b)} Jika $f$ merupakan fungsi nonnegatif yang
$\mu_H$-terintegralkan, terdapat barisan tak-menurun
$\sequencen{f_n}$ dari fungsi-fungsi $\mu_H$-sederhana nonnegatif
yang konvergen ke $f\,\,\mu_H$-hampir di mana-mana; sekarang
$\sequencen{\tilde f_n}$ merupakan barisan tak-menurun dari
fungsi-fungsi $\mu$-sederhana yang konvergen ke
$\tilde f\,\,\mu$-hampir di mana-mana (131Ca), dan

\Centerline{$ \sup_{n\in\Bbb N}\int \tilde f_nd\mu
=\sup_{n\in\Bbb N}\int f_nd\mu_H
=\int fd\mu_H<\infty$,}

\noindent sehingga $\int\tilde fd\mu$ ada dan sama dengan
$\int fd\mu_H$.

\medskip

{\bf (c)} Jika $f$ $\mu_H$-terintegralkan, fungsi itu dapat dinyatakan
sebagai $f_1-f_2$, dengan $f_1$ dan $f_2$ fungsi-fungsi nonnegatif
yang $\mu_H$-terintegralkan, sehingga
$\tilde f=\tilde f_1-\tilde f_2$ dan

\Centerline{$ \int\tilde fd\mu=\int\tilde f_1d\mu-\int\tilde f_2d\mu=\int
f_1d\mu_H-\int f_2d\mu_H=\int fd\mu_H$.}

\medskip

{\bf (d)} Sekarang andaikan bahwa $\tilde f$ $\mu$-terintegralkan.
Dalam hal ini terdapat himpunan $\mu$-koterabaikan $E\in\Sigma$
sedemikian sehingga $E\subseteq\dom\tilde f$ dan $\tilde f\restr E$
bersifat $\Sigma$-terukur (122P). Tentu saja $\mu(H\setminus E)=0$,
sehingga $E\cap H$ bersifat $\mu_H$-koterabaikan; selain itu, untuk
setiap $a\in\Bbb R$,

\Centerline{$\{x:x\in E\cap H,\,f(x)\ge a\}
=H\cap\{x:x\in E,\,\tilde f(x)\ge a\}\in\Sigma_H$,}

\noindent sehingga $f\restr(E\cap H)$ bersifat $\Sigma_H$-terukur,
dan $f$ terukur secara virtual terhadap $\mu_H$ serta didefinisikan
$\mu_H$-hampir di mana-mana. Selanjutnya, untuk $\epsilon>0$,

\Centerline{$\mu_H\{x:x\in E\cap H,\,|f(x)|\ge\epsilon\}
=\mu\{x:x\in E,\,|\tilde f(x)|\ge\epsilon\}<\infty$,}

\noindent sedangkan jika $g$ suatu fungsi $\mu_H$-sederhana dan
$g\le|f|\,\,\mu_H$-hampir di mana-mana, maka
$\tilde g\le|\tilde f|\,\,\mu$-hampir di mana-mana dan

\Centerline{$\int g\,d\mu_H=\int\tilde g\,d\mu
\le\int|\tilde  f|d\mu<\infty$.}

\noindent Menurut kriteria dalam 122J dan 122P, $f$
$\mu_H$-terintegralkan, sehingga sekali lagi kita memperoleh
$\int fd\mu_H=\int\tilde fd\mu$.
}%end of proof of 131E


\leader{131F}{Korolari} Misalkan $(X,\Sigma,\mu)$ suatu ruang ukur dan
$f$ suatu fungsi bernilai real yang didefinisikan pada suatu
subhimpunan $\dom f$ dari $X$.

(a) Jika $H\in\Sigma$ dan $f$ didefinisikan hampir di mana-mana dalam
$X$, maka $f\restr H$ $\mu_H$-terintegralkan jika dan hanya jika
$f\times\chi H$ $\mu$-terintegralkan, dan dalam hal ini
$\int_Hf=\int f\times\chi H$.

(b) Jika $f$ $\mu$-terintegralkan, maka $f\ge 0$ a.e.\ jika dan hanya
jika $\int_Hf\ge 0$ untuk setiap $H\in\Sigma$.

(c) Jika $f$ $\mu$-terintegralkan, maka $f=0$ a.e.\ jika dan hanya
jika $\int_Hf=0$ untuk setiap $H\in\Sigma$.

\proof{{\bf (a)} Karena $\dom f$ bersifat $\mu$-koterabaikan,
$(f\restr H)\ssptilde$, sebagaimana didefinisikan dalam 131E, sama
dengan $f\times\chi H\,\,\mu$-hampir di mana-mana, sehingga, menurut
131E,

\Centerline{$ \int_Hfd\mu
=\int(f\restr H)\ssptilde d\mu
=\int(f\times\chi H)d\mu$}

\noindent jika salah satu integral itu ada; dalam hal ini integral
lainnya juga ada.

\medskip

{\bf (b)(i)} Jika $f\ge 0\,\,\mu$-hampir di mana-mana, maka untuk
setiap $H\in\Sigma$ kita harus mempunyai
$f\restr H\ge 0\,\,\mu_H$-hampir di mana-mana, sehingga
$\int_Hf=\int(f\restr H)d\mu_H\ge 0$.

\medskip

\quad{\bf (ii)} Jika $\int_Hf\ge 0$ untuk setiap $H\in\Sigma$,
misalkan $E\in\Sigma$ suatu subhimpunan koterabaikan dari $\dom f$
sedemikian sehingga $f\restr E$ terukur.
Tetapkan $F=\{x:x\in E,\,f(x)<0\}$. Maka $\int_Ff\ge 0$; menurut
122Rc, diperoleh $f\restr F=0\,\,\mu_F$-hampir di mana-mana, dan ini
hanya mungkin jika $\mu F=0$; dalam hal itu
$f\ge 0\,\,\mu$-hampir di mana-mana.

\medskip

{\bf (c)} Terapkan (b) pada $f$ dan pada $-f$ untuk melihat bahwa
$f\le 0\le f$ hampir di mana-mana.
}%end of proof of 131F

\leader{131G}{Korolari} Misalkan $(X,\Sigma,\mu)$ suatu ruang ukur dan
$H\in\Sigma$ suatu himpunan koterabaikan. Jika $f$ sembarang fungsi
bernilai real yang didefinisikan pada suatu subhimpunan dari $X$, maka,
jika salah satu ruas dalam $\int_Hf=\int f$ terdefinisi, ruas lainnya
juga demikian dan persamaan tersebut berlaku.

\proof{ Dalam bahasa 131E, $f=(f\restr H)\ssptilde\,\,\mu$-hampir di
mana-mana, sehingga

\Centerline{$\int f=\int(f\restr H)\ssptilde=\int_Hf$}

\noindent jika salah satu integral itu terdefinisi; dalam hal ini
integral lainnya juga terdefinisi.
}%end of proof of 131G

\leader{131H}{Korolari} Misalkan $(X,\Sigma,\mu)$ suatu ruang ukur dan
$f$, $g$ dua fungsi bernilai real yang $\mu$-terintegralkan.

(a) Jika $\int_Hf\ge\int_Hg$ untuk setiap $H\in\Sigma$, maka
$f\ge g$ a.e.

(b) Jika $\int_Hf=\int_Hg$ untuk setiap $H\in\Sigma$, maka
$f=g$ a.e.

\proof{ Terapkan 131Fb-131Fc pada $f-g$.
}%end of proof of 131H

\exercises{
\leader{131X}{Latihan dasar $\pmb{>}$(a)}
%\spheader 131Xa
Misalkan $(X,\Sigma,\mu)$ suatu ruang ukur, dan $f$ suatu fungsi
bernilai real yang terintegralkan pada $X$. Untuk $E\in\Sigma$
tetapkan $\nu E=\int_Ef$. (i) Tunjukkan bahwa jika $E$, $F$ merupakan
anggota-anggota saling lepas dari $\Sigma$, maka
$\nu(E\cup F)=\nu E+\nu F$. \Hint{131E.}
(ii) Tunjukkan bahwa jika $\sequencen{E_n}$ suatu barisan saling lepas
di dalam $\Sigma$, maka
$\nu(\bigcup_{n\in\Bbb N}E_n)=\sum_{n=0}^{\infty}\nu E_n$.
\Hint{123C.} (iii) Tunjukkan bahwa jika $f$ nonnegatif, maka
$(X,\Sigma,\nu)$ merupakan suatu ruang ukur.

\sqheader 131Xb Misalkan $\mu$ ukuran Lebesgue pada $\Bbb R$.
(i) Tunjukkan bahwa setiap kali $a\le b$ dalam $\Bbb R$ dan $f$ suatu
fungsi bernilai real dengan $\dom f\subseteq\Bbb R$, maka

\Centerline{$\int_{\ooint{a,b}}fd\mu=\int_{\coint{a,b}}fd\mu
=\int_{\ocint{a,b}}fd\mu=\int_{[a,b]}fd\mu$}

\noindent jika salah satunya terdefinisi. \Hint{terapkan 131E pada
empat versi berbeda dari $\tilde f$.} Tuliskan $\int_a^bfd\mu$ untuk
nilai bersama ini.
(ii) Tunjukkan bahwa jika $a\le b\le c$ dalam $\Bbb R$, maka, untuk
setiap fungsi bernilai real $f$,
$\int_a^cfd\mu=\int_a^bfd\mu+\int_b^cfd\mu$ jika salah satu ruas
terdefinisi; dalam hal ini ruas lainnya juga terdefinisi. (iii)
Tunjukkan bahwa jika $f$ terintegralkan pada
$\Bbb R$, maka $(a,b)\mapsto\int_a^bfd\mu$ kontinu.
\Hint{{\it Salah satu\/}: tinjau terlebih dahulu fungsi-fungsi
sederhana $f$, {\it atau\/} tinjau
$\lim_{n\to\infty}\int_{a_n}^bfd\mu$ untuk barisan-barisan monoton
$\sequencen{a_n}$.}

\spheader 131Xc Misalkan $g:\Bbb R\to\Bbb R$ suatu fungsi
tak-menurun dan $\mu_g$ ukuran Lebesgue-Stieltjes yang terkait
(114Xa). (i) Tunjukkan bahwa jika $a\le b\le c$ dalam $\Bbb R$, maka,
untuk setiap fungsi bernilai real $f$,
$\int_{\coint{a,c}}fd\mu_g
=\int_{\coint{a,b}}fd\mu_g+\int_{\coint{b,c}}fd\mu_g$ jika salah satu
ruas terdefinisi; dalam hal ini ruas lainnya juga terdefinisi. (ii)
Tunjukkan bahwa jika $f$ $\mu_g$-terintegralkan
pada $\Bbb R$, maka
$(a,b)\mapsto\int_{\coint{a,b}}fd\mu_g$ kontinu pada
$\{(a,b):a\le b,\,g$ kontinu di $a$ maupun $b\}$.

\leader{131Y}{Latihan lanjutan (a)}
%\spheader 131Ya
Misalkan $(X,\Sigma,\mu)$ suatu ruang ukur dan $E\in\Sigma$ suatu
himpunan terukur dengan ukuran berhingga. Misalkan $\sequencen{f_n}$
suatu barisan fungsi bernilai real yang terukur, dengan domain-domain
terukur\footnote{Saya berterima kasih kepada P. Wallace Thompson karena
telah menunjukkan bahwa klausa ini, atau sesuatu yang setara dengannya,
diperlukan.}, sedemikian sehingga
$f=\lim_{n\to\infty}f_n$ didefinisikan hampir di mana-mana dalam $E$
(sesuai dengan kesepakatan-kesepakatan 121Fa). Tunjukkan bahwa untuk
setiap $\epsilon>0$ terdapat himpunan terukur $F\subseteq E$ sedemikian
sehingga $\mu(E\setminus F)\le\epsilon$ dan $\sequencen{f_n}$
konvergen seragam ke $f$ pada $F$. (Ini adalah {\bf teorema Egorov}.)

}%end of exercises

\endnotes{
\Notesheader{131} Jika Anda menginginkan definisi singkat bagi
$\int_Hf$ untuk $H$ terukur, definisi dalam 131E tampaknya yang paling
sederhana, karena memungkinkan Anda menghindari konsep `ukuran
subruang' sepenuhnya. Namun, menurut saya kita sungguh perlu dapat
berbicara, misalnya, tentang `ukuran Lebesgue pada $[0,1]$', dengan
memaksudkan ukuran subruang $\mu_{[0,1]}$ ketika $\mu$ adalah ukuran
Lebesgue pada $\Bbb R$.

Bagian ini memuat cukup banyak analisis teknis yang terperinci.
Masalahnya ialah bahwa mulai dari 131A, pada umumnya terdapat
sekurang-kurangnya dua ukuran yang berperan, dan bahasa biasa dalam
teori ukuran---kata-kata seperti `terukur' dan `terintegralkan'---
menjadi tidak dapat dipercaya dalam konteks semacam itu, karena bahasa
tersebut menghilangkan keterangan penting mengenai aljabar-$\sigma$
atau ukuran mana yang sedang ditinjau. Karena itu saya harus
menggunakan istilah yang kurang anggun tetapi lebih eksplisit. Namun,
saya berharap bahwa setelah Anda menelaah dengan saksama hasil-hasil
seperti 131F, Anda akan merasa bahwa pola yang terbentuk cukup koheren.
Aturan umumnya ialah bahwa untuk subruang-subruang {\it terukur} tidak
ada kejutan yang serius (131Cb, 131Fa).

Perlu saya catat bahwa terdapat pula definisi standar bagi ukuran
subruang pada subhimpunan-subhimpunan {\it tak-terukur} dari suatu ruang
ukur. Definisi itu telah saya berikan dalam 113Yb; untuk teori
integrasi yang memperluas hasil-hasil di atas, saya akan menunggu hingga
\S214. Terdapat kesulitan-kesulitan tambahan yang berarti, dan
keumuman tambahan tersebut tidak sering diperlukan dalam penerapan
elementer.

Saya hendak menarik perhatian Anda pada 131Fb-131Fc dan 131Xa-131Xc;
semuanya merupakan langkah awal untuk memahami `integral tak tentu',
yaitu fungsional-fungsional
$E\mapsto\int_Ef:\Sigma\to\Bbb R$ ketika $f$ suatu fungsi
terintegralkan. Saya akan kembali membahasnya dalam Bab 22 dan 23.
}%end of notes



\discrpage






